% ============================================================
%  TechSource - Local LLM + MATLAB Agentic Toolkit Setup Guide
%  LM Studio (Bionic) + MATLAB MCP Server
%  + MATLAB Agentic Toolkit + Simulink Agentic Toolkit
% ============================================================
\documentclass[11pt, letterpaper]{article}

\usepackage[top=0.7in, bottom=0.9in, left=0.85in, right=0.85in, footskip=0.5in]{geometry}
\usepackage{graphicx}
\usepackage{xcolor}
\usepackage[T1]{fontenc}
\usepackage{textcomp}
\usepackage{helvet}
\usepackage{microtype}
\usepackage[nobottomtitles*]{titlesec}
\usepackage{titling}
\usepackage{enumitem}
\usepackage{hyperref}
\usepackage{fancyhdr}
\usepackage{booktabs}
\usepackage{array}
\usepackage{parskip}
\usepackage{tcolorbox}
\usepackage{caption}
\usepackage{float}
\usepackage{amssymb}

\renewcommand{\familydefault}{\sfdefault}

\definecolor{TSorange}{RGB}{230, 126, 14}
\definecolor{TSgray}{RGB}{100, 100, 100}
\definecolor{TSdark}{RGB}{30, 30, 30}
\definecolor{TSlightgray}{RGB}{245, 245, 245}
\definecolor{TSblue}{RGB}{30, 90, 160}
\definecolor{TSboxbg}{RGB}{255, 248, 235}
\definecolor{TScodebg}{RGB}{243, 245, 248}
\definecolor{TSframe}{RGB}{185, 185, 185}

\hypersetup{
  colorlinks = true,
  linkcolor  = TSorange,
  urlcolor   = TSblue,
  citecolor  = TSgray,
  pdftitle   = {Local LLM + MATLAB Agentic Toolkit Setup},
  pdfauthor  = {TechSource},
}

\titleformat{\section}
  {\color{TSorange}\large\bfseries\sffamily}
  {Step \thesection.}{0.6em}{}
  [{\color{TSorange}\titlerule[0.8pt]}]

\titleformat{\subsection}
  {\color{TSgray}\normalsize\bfseries\sffamily}
  {\thesubsection}{0.6em}{}

\titlespacing{\section}{0pt}{10pt}{5pt}
\titlespacing{\subsection}{0pt}{7pt}{3pt}

% Never leave a step heading stranded at the foot of a page
\renewcommand{\bottomtitlespace}{0.10\textheight}

\pagestyle{fancy}
\fancyhf{}
\setlength{\headheight}{46pt}
\fancyhead[L]{\includegraphics[height=0.42in]{techsource_logo.png}}
\fancyhead[R]{\includegraphics[height=0.42in]{mw-reseller-badge-rgb-def.png}}
\fancyfoot[C]{\color{TSgray}\scriptsize\thepage}
\renewcommand{\headrulewidth}{0.4pt}

\fancypagestyle{plain}{\fancyhf{}\renewcommand{\headrulewidth}{0pt}\renewcommand{\footrulewidth}{0pt}}

\captionsetup{
  font={small,color=TSgray},
  labelfont={bf,color=TSorange},
  format=plain,
  justification=centering,
  skip=4pt,
}

\tcbuselibrary{skins, breakable}

\newtcolorbox{notebox}{
  colback=TSboxbg, colframe=TSorange,
  fonttitle=\bfseries\sffamily\small,
  title={$\bigstar$\; Note},
  left=6pt, right=6pt, top=4pt, bottom=4pt, breakable,
  fontupper=\raggedright,
}

\newtcolorbox{stepbox}[1]{
  colback=TSlightgray, colframe=TSgray,
  fonttitle=\bfseries\sffamily\small\color{TSorange},
  title={#1},
  left=6pt, right=6pt, top=4pt, bottom=4pt, breakable, arc=3pt,
  fontupper=\raggedright,
}

% Command / path box -- for anything the reader must type or paste
\newtcolorbox{codebox}{
  colback=TScodebg, colframe=TSframe,
  boxrule=0.5pt, arc=2pt,
  left=6pt, right=6pt, top=5pt, bottom=5pt, breakable,
  fontupper=\ttfamily\small\raggedright,
}

% ------------------------------------------------------------
%  Uniform screenshot command
%  Usage: \screenshot[label]{filename}{Caption}
%
%  Every screenshot is centred inside a frame of EXACTLY the same
%  width and height, so all figures line up down the page no matter
%  what the source image resolution is. To resize every screenshot
%  in the document, edit \shotwidth below -- nothing else.
%
%  HOW TO INSERT A SCREENSHOT:
%    1. Save your PNG/JPG in the same folder as this .tex file
%    2. Add a \screenshot{file.png}{Caption} line where you want it
%    3. Recompile twice: pdflatex local_llm_setup_instructions.tex
% ------------------------------------------------------------
\newlength{\shotwidth}
\newlength{\shotheight}
\setlength{\shotwidth}{0.73\textwidth}
\setlength{\shotheight}{0.538\shotwidth}

% Slightly tighter, more flexible spacing around the framed figures so a
% step and its two screenshots stay together on one page.
\setlength{\intextsep}{10pt plus 3pt minus 3pt}

\newcommand{\screenshot}[3][]{%
  \begin{figure}[H]
    \centering
    {\setlength{\fboxsep}{0pt}\setlength{\fboxrule}{0.6pt}%
     \fcolorbox{TSframe}{white}{%
       \begin{minipage}[c][\shotheight][c]{\shotwidth}%
         \setlength{\parskip}{0pt}%
         \centering
         \includegraphics[width=\shotwidth, height=\shotheight, keepaspectratio]{#2}%
       \end{minipage}}}%
    \ifx&#3&\else\caption{#3}\fi
    \ifx&#1&\else\label{#1}\fi
  \end{figure}%
}

\setlist[enumerate,1]{label=\textbf{\arabic*.}, leftmargin=*, itemsep=4pt, topsep=4pt}
\setlist[itemize,1]{leftmargin=*, itemsep=3pt, topsep=3pt, label=\textcolor{TSorange}{\textbullet}}

\setlength{\parskip}{5pt}
\setlength{\parindent}{0pt}

\newcommand{\tick}{\textcolor{TSorange}{$\square$}}
\newcommand{\arrow}{\,$\rightarrow$\,}

% ============================================================
\begin{document}
% ============================================================

% -- Compact Title Banner -------------------------------------
\begin{center}
  \color{TSorange}\rule{\textwidth}{2pt}\\[5pt]
  {\fontsize{18}{22}\selectfont\bfseries\color{TSdark}
    Local LLM + MATLAB Agentic Toolkit Setup Guide}\\[4pt]
  {\small\color{TSgray}
    LM Studio \textendash{} MATLAB MCP Server \textendash{} MATLAB \& Simulink Agentic Toolkits}\\[2pt]
  {\small\color{TSgray}
    Platform: Windows \quad\textbar\quad MATLAB R2026a
    \quad\textbar\quad \textbf{Runs Fully Offline}
    \quad\textbar\quad Revised 19 August 2026}\\[5pt]
  \color{TSorange}\rule{\textwidth}{2pt}
\end{center}

\vspace{6pt}

% ============================================================
\section*{Overview}
% ============================================================

This guide sets up a \textbf{fully local AI assistant for MATLAB and Simulink}. The language model
runs on your own machine, so your code and models never leave your computer.

\begin{notebox}
  \textbf{Once the setup is done, the whole thing runs offline.} The model, the MCP server, the
  tools and the skills all live on your own disk --- nothing is sent to a cloud service, and no
  internet connection is needed to use it. You only need to be online for the \textbf{downloads}
  in Steps 1, 3, 4, 5 and 10.
\end{notebox}

Four pieces work together:

\begin{center}
\small
\begin{tabular}{@{}>{\raggedright\arraybackslash}p{0.30\textwidth} >{\raggedright\arraybackslash}p{0.62\textwidth}@{}}
\toprule
\textbf{Component} & \textbf{What it does} \\
\midrule
LM Studio (Bionic)          & Runs the language model locally and hosts the chat / agent interface. \\
MATLAB MCP Server           & Lets the agent drive MATLAB, starting it if it is not already open. \\
MATLAB Agentic Toolkit      & Adds MATLAB tools and skills the agent can use. \\
Simulink Agentic Toolkit    & Adds Simulink, Stateflow and System Composer tools and skills. \\
\bottomrule
\end{tabular}
\end{center}

\begin{notebox}
  Set aside about \textbf{1--2 hours}, most of which is downloading. You will need roughly
  \textbf{15 GB of free disk space} for LM Studio, one model, and the toolkits.
\end{notebox}

\textbf{What you will need:}
\begin{itemize}
  \item A \textbf{Windows} PC where you are allowed to install software
  \item \textbf{MATLAB R2026a} installed locally, plus \textbf{Simulink} for the Simulink toolkit
  \item A stable internet connection --- for the setup downloads only
  \item A GPU with 8 GB or more of VRAM is recommended, but not required
\end{itemize}

\begin{notebox}
  Please complete this setup \textbf{at least one day before} the workshop. If you get stuck,
  contact \texttt{matlabday\_support@techsource-asia.com} with a screenshot of the step that
  failed.
\end{notebox}

\newpage

% ============================================================
\section{Install LM Studio}
% ============================================================

\begin{enumerate}
  \item Go to \url{https://lmstudio.ai/download} and click
        \textbf{Download LM Studio Bionic}. The installer is about 600 MB.

  \screenshot[fig:lms-download]{lm_studio_installation.png}{Downloading the LM Studio Bionic installer}

  \item Run the downloaded \texttt{Bionic-\textit{version}-x64.exe} and accept the default options.

  \screenshot[fig:lms-install]{lm_studio.png}{The LM Studio (Bionic) installer running}
\end{enumerate}

\begin{notebox}
  \textbf{Note down the installation path} shown by the installer. Separately, LM Studio keeps
  your per-user data --- models, skills and settings --- in a hidden folder in your user profile:
  \begin{codebox}
  \path{C:\Users\<YourUserName>\.lmstudio}
  \end{codebox}
  You will need that second path in \textbf{Step 9}.
\end{notebox}

% ============================================================
\section{Open the LM Studio Settings}
% ============================================================

\begin{enumerate}
  \item Launch LM Studio. On a fresh install you land on the \textbf{Create your first project}
        screen. Click the \textbf{sidebar toggle} in the top-left corner.

  \screenshot[fig:lms-sidebar]{lm_studio_1.png}{Click the sidebar toggle at the top-left}

  \item Click \textbf{Settings} at the bottom of the sidebar.

  \screenshot[fig:lms-settings]{2.png}{Settings sits at the bottom of the expanded sidebar}
\end{enumerate}

% ============================================================
\section{Choose and Download a Model}
% ============================================================

\begin{enumerate}
  \item In the Settings sidebar, under \textbf{Local Models}, click \textbf{Explore}.
  \item Pick a model from the staff picks, or search Hugging Face. This guide uses
        \textbf{Gemma 4 12B QAT} (7.15 GB), which fits comfortably on most workshop machines.
  \item Click \textbf{Download} and let it finish before continuing.

  \screenshot[fig:lms-model]{3.png}{Selecting Gemma 4 12B QAT in Explore and starting the download}
\end{enumerate}

\begin{notebox}
  Check the \textbf{Capabilities} row in the model details before downloading. The model
  \textbf{must list \texttt{Tools}} --- without it the model cannot call the MATLAB MCP tools and
  nothing later in this guide will work. \texttt{Vision} and \texttt{Reasoning} are useful extras.
\end{notebox}

% ============================================================
\section{Download the Agentic Toolkit Installer}
% ============================================================

\begin{enumerate}
  \item Go to the Simulink Agentic Toolkit repository:\\
        \url{https://github.com/matlab/simulink-agentic-toolkit}
  \item In the right-hand column under \textbf{Releases}, click the release marked
        \textbf{Latest}.

  \screenshot[fig:gh-release]{5.png}{Open the latest release from the repository page}

  \item Scroll down to \textbf{Assets} and download \texttt{agenticToolkitInstaller.mltbx}.

  \screenshot[fig:gh-asset]{6.png}{Download agenticToolkitInstaller.mltbx from the release assets}
\end{enumerate}

\begin{notebox}
  \texttt{agenticToolkitInstaller.mltbx} is the \textbf{only} file you need. It is a small
  installer, about 76 KB, that fetches the MCP server and both toolkits for you. Do \textbf{not}
  download \texttt{simulink-agentic-toolkit.mltbx} separately.
\end{notebox}

% ============================================================
\section{Install the Toolkits in MATLAB}
% ============================================================

\begin{enumerate}
  \item Open MATLAB and set the \textbf{Current Folder} to the folder you downloaded the
        installer into, for example \texttt{Downloads}.
  \item In the \textbf{Command Window}, run the following, substituting your own path:

  \begin{codebox}
  \path{matlab.internal.open.openmltbx("C:\Users\TSA\Downloads\agenticToolkitInstaller.mltbx")}
  \end{codebox}

  \screenshot[fig:ml-openmltbx]{7.png}{Running openmltbx on the downloaded .mltbx file}

  \item When \textbf{Installation Complete} appears, click \textbf{Close}.

  \screenshot[fig:ml-installed]{8.png}{The installer confirms the toolbox is installed}

  \item Now run the setup function:

  \begin{codebox}
  setupAgenticToolkit("install", DisableTelemetry=true)
  \end{codebox}

  \item At \textbf{Enter selection}, type \texttt{3} for \textbf{Both} and press Enter. Review the
        plan, then type \texttt{y} at \textbf{Proceed? [Y/n]} and press Enter.

  \screenshot[fig:ml-setup]{9.png}{Select option 3 (Both), then confirm with y}

  \item The setup downloads and installs the MCP server and both toolkits. When it asks
        \textbf{``Would you like to configure your coding agent now?''}, type \texttt{n}
        and press Enter.

  \screenshot[fig:ml-agent-n]{10.png}{Answer n --- LM Studio is registered by hand in Step 7}
\end{enumerate}

\begin{notebox}
  Answer \texttt{n} because the automatic configuration only covers Claude Code, Codex, GitHub
  Copilot, Amp and Gemini CLI. LM Studio is registered manually in \textbf{Step 7}.
\end{notebox}

% ============================================================
\section{Copy the MATLAB MCP Server Path}
% ============================================================

\begin{enumerate}
  \item In the setup output, find the line beginning
        \textbf{``MCP server binary installed to:''} and select and copy the full path.

  \screenshot[fig:ml-mcp-path]{11.png}{Copy the MCP server binary path from the install log}
\end{enumerate}

On a default installation the path looks like this:

\begin{codebox}
\path{C:\Users\TSA\.matlab\agentic-toolkits\bin\matlab-mcp-server.exe}
\end{codebox}

\begin{notebox}
  Leave the MATLAB Command Window open and do not clear it. You will come back for a second
  path in the middle of the next step.
\end{notebox}

% ============================================================
\section{Register the MCP Server in LM Studio}
% ============================================================

\begin{enumerate}
  \item In LM Studio, go to \textbf{Settings}\arrow\textbf{Connected Apps}. At the bottom, under
        \textbf{Manual MCP server setup}, click \textbf{+ Add custom MCP}.

  \screenshot[fig:lms-connected]{4.png}{Add custom MCP under Manual MCP server setup}

  \item Fill in the form:
  \begin{itemize}
    \item \textbf{Name}: \texttt{MATLAB MCP Core Server}
    \item \textbf{Connection}: \texttt{On this computer}
    \item \textbf{Command}: paste the MCP server path you copied in Step 6
  \end{itemize}

  \item Switch back to MATLAB and copy the \textbf{Simulink Agentic Toolkit} path from the line
        \textbf{``Simulink Agentic Toolkit installed to:''}.

  \screenshot[fig:ml-slk-path]{12.png}{Copy the Simulink Agentic Toolkit installation path}

  \item Back in LM Studio, turn on \textbf{Show advanced options}. Under \textbf{Arguments},
        add one argument built from the path you just copied:

  \begin{codebox}
  \path{--extension-file=C:\Users\TSA\.matlab\agentic-toolkits\simulink\tools\tools.json}
  \end{codebox}

  \begin{sloppypar}
  Replace \path{C:\Users\TSA\.matlab\agentic-toolkits\simulink} with your own Simulink Agentic
  Toolkit path, and keep the trailing \path{\tools\tools.json} exactly as shown.
  \end{sloppypar}

  \screenshot[fig:lms-mcp-args]{13.png}{The Command field and the extension-file argument}

  \item Scroll to the bottom of the form and click \textbf{Add MCP}.

  \screenshot[fig:lms-add-mcp]{14.png}{Click Add MCP to save the server}

  \item Back on the \textbf{Connected Apps} page, confirm the server reports
        \textbf{Connected \textbullet{} 13 tools ready}.

  \screenshot[fig:lms-connected-ok]{15.png}{The MATLAB MCP Core Server connected with 13 tools}
\end{enumerate}

\begin{notebox}
  If you do not see \textbf{13 tools ready}, check these in order:
  \begin{itemize}
    \item Toggle the server \textbf{off and on again} --- LM Studio only launches the server when
          it is enabled.
    \item The \textbf{Command} path points at \texttt{matlab-mcp-server.exe} and that file exists.
    \item The \texttt{--extension-file} argument ends in \path{\tools\tools.json}
          and that file exists. Without this argument you will see fewer tools, because the
          Simulink tools are never loaded.
  \end{itemize}
\end{notebox}

% ============================================================
\section{Increase the Context Length}
% ============================================================

\begin{enumerate}
  \item In Settings, under \textbf{Local Models}, click \textbf{Local Model Defaults}.
  \item Set \textbf{Minimum AutoFit context length} to \texttt{325000}.

  \screenshot[fig:lms-context]{16.png}{Raise Minimum AutoFit context length to 325000}
\end{enumerate}

\begin{notebox}
  The toolkit tool definitions and skills are large. With too little context the agent silently
  truncates them and tool calls start to fail. If a model then \textbf{refuses to load}, your
  machine cannot provide that much context --- lower this value or choose a smaller model.
\end{notebox}

% ============================================================
\section{Open the LM Studio Skills Folder}
% ============================================================

\begin{enumerate}
  \item Open File Explorer, paste the following into the address bar, and press Enter:

  \begin{codebox}
  \path{%USERPROFILE%\.lmstudio\skills}
  \end{codebox}

  \screenshot[fig:skills-folder]{17.png}{The hidden .lmstudio skills folder, empty on a fresh install}

  \item Leave this window open --- you will paste into it in the next step.
\end{enumerate}

\begin{notebox}
  \texttt{.lmstudio} is hidden, so you will not find it by browsing. Paste the path above, or turn
  on \textbf{View}\arrow\textbf{Show}\arrow\textbf{Hidden items}. If the \texttt{skills} folder
  does not exist yet, create it.
\end{notebox}

% ============================================================
\section{Install the MATLAB and Simulink Skills}
% ============================================================

\begin{enumerate}
  \item Open both toolkit repositories in your browser:
  \begin{itemize}
    \item \url{https://github.com/matlab/matlab-agentic-toolkit/}
    \item \url{https://github.com/matlab/simulink-agentic-toolkit/}
  \end{itemize}

  \item On each repository, click the green \textbf{Code} button and choose
        \textbf{Download ZIP}.

  \screenshot[fig:gh-zip]{18.png}{Code, then Download ZIP, on each repository}

  \item Extract both ZIP files. Open
        \path{matlab-agentic-toolkit-main\skills-catalog}, select
        \textbf{all of the folders}, and copy them.

  \screenshot[fig:matlab-skills]{19.png}{The MATLAB skills catalog --- copy all folders}

  \item Paste them into the \path{.lmstudio\skills} folder from Step 9.

  \item Repeat for the Simulink toolkit: open
        \path{simulink-agentic-toolkit-main\skills-catalog}, copy all of the
        folders, and paste them into the same \texttt{skills} folder.

  \screenshot[fig:simulink-skills]{20.png}{The Simulink skills catalog --- copy all folders}
\end{enumerate}

\begin{notebox}
  Copy the \textbf{contents} of \texttt{skills-catalog}, not the \texttt{skills-catalog} folder
  itself. Each skill category must sit directly inside
  \path{.lmstudio\skills}. You can skip \texttt{README.md}.
\end{notebox}

% ============================================================
\section{Verify the Skills Are Loaded}
% ============================================================

\begin{enumerate}
  \item In LM Studio, go to \textbf{Settings}\arrow\textbf{Skills}.
  \item Confirm the \textbf{Bionic Skills} list is populated, with entries such as
        \texttt{building-simulink-models}, \texttt{checking-model-compliance} and
        \texttt{authoring-simulink-inputs}.

  \screenshot[fig:lms-skills]{21.png}{The Bionic Skills list after copying both catalogs}
\end{enumerate}

\begin{notebox}
  If the list is empty, restart LM Studio. If it is still empty, check the folder layout --- each
  skill folder must contain a \texttt{SKILL.md} file.
\end{notebox}

% ============================================================
\section{Run Your First Prompt}
% ============================================================

Your setup is complete. You do not need to start MATLAB yourself --- on the first tool call the
MCP server launches MATLAB for you. If MATLAB is already open, the agent attaches to that session
instead.

\begin{enumerate}
  \item Click \textbf{Back to app} to leave Settings, then click \textbf{New Project}. Give the
        project a name, leave \textbf{Allow coding} switched on, pick a folder for the
        \textbf{Local workspace}, and click \textbf{Create}.

  \screenshot[fig:first-project]{22.png}{Creating a project with Allow coding switched on}

  \item In the prompt box, click the model selector and choose the model you downloaded in
        Step 3.

  \screenshot[fig:first-model]{23.png}{Choosing the model in the prompt box}

  \item Send this prompt:

  \begin{codebox}
  What version of MATLAB is running? List the installed toolboxes.
  \end{codebox}

  \item Watch the session. The agent shows each tool as it runs --- the first is
        \textbf{Called MATLAB MCP Core Server}, and MATLAB starts up alongside it.

  \screenshot[fig:first-prompt]{24.png}{The first tool call, with MATLAB launching automatically}

  \item Wait for the answer. It should report \textbf{your} MATLAB release and \textbf{your}
        toolbox list. Expand \textbf{Called 1 tool} and \textbf{Reasoned} to see which tool the
        model picked --- here it chose \texttt{detect\_matlab\_toolboxes}.

  \screenshot[fig:first-result]{25.png}{The reply: release and toolbox list read live from MATLAB}
\end{enumerate}

\begin{notebox}
  The \textbf{first} prompt is slow, because MATLAB has to start before the tool can return.
  Later prompts in the same session reuse the running MATLAB and respond much faster.
\end{notebox}

\begin{notebox}
  \textbf{Everything you just ran was offline.} To prove it, disconnect from the network and send
  the prompt again --- the model, the MCP server, the tools and the skills are all on your
  machine. The one exception is MATLAB's own licence check, if yours is a network licence.
\end{notebox}

% ------------------------------------------------------------
%  Add further exercise prompts below, following the same
%  pattern. Every screenshot is framed at the same width as the
%  ones earlier in the guide automatically -- there is nothing
%  else to adjust.
%
%  \screenshot[fig:prompt-slk]{25.png}{The agent building a Simulink model}
% ------------------------------------------------------------

\newpage

% ============================================================
%  Closing reference page -- checklist and troubleshooting are
%  kept together on their own page so it can be printed alone.
% ============================================================
\section*{Setup Checklist}
% ============================================================

Run through this list before the workshop:

\begin{itemize}[label={\tick}]
  \item LM Studio is installed and opens without errors
  \item A model with the \texttt{Tools} capability has finished downloading
  \item MATLAB R2026a is installed locally, with Simulink
  \item \texttt{setupAgenticToolkit} completed and reported \textbf{Artifacts installed}
  \item \textbf{Connected Apps} shows \textbf{Connected \textbullet{} 13 tools ready}
  \item \textbf{Minimum AutoFit context length} is set to \texttt{325000}
  \item \textbf{Settings}\arrow\textbf{Skills} lists the MATLAB and Simulink skills
  \item The first prompt returned \textbf{your} MATLAB release and toolbox list
\end{itemize}

% ============================================================
\section*{Troubleshooting}
% ============================================================

\begin{center}
\small
\begin{tabular}{@{}>{\raggedright\arraybackslash}p{0.38\textwidth} >{\raggedright\arraybackslash}p{0.54\textwidth}@{}}
\toprule
\textbf{Symptom} & \textbf{Most likely cause} \\
\midrule
MCP server will not connect
  & The \textbf{Command} path is wrong, or \texttt{matlab-mcp-server.exe} is missing. \\
Connected, but fewer than 13 tools
  & The \texttt{--extension-file} argument is missing, or points at a \texttt{tools.json} that does not exist. \\
Model fails to load
  & Not enough RAM or VRAM for a 325{,}000-token context. Lower the value, or use a smaller model. \\
Answer arrives with no \textbf{Called} row
  & The model is answering from general knowledge, not from MATLAB. Check that it has the
    \texttt{Tools} capability, and that Step 7 reads 13 tools ready. \\
Skills list is empty
  & The skills were copied one level too deep, or LM Studio needs a restart. \\
\texttt{Unrecognized function or variable}
  & \texttt{setupAgenticToolkit} only exists after the \texttt{.mltbx} is installed --- redo Step 5. \\
\bottomrule
\end{tabular}
\end{center}

\vfill
\begin{center}
  {\color{TSgray}\scriptsize\textit{\copyright{} \the\year{} TechSource. All rights reserved.}}
\end{center}

\end{document}
